\documentclass[aspectratio=169]{beamer}

\usetheme{avalanche}

\hypersetup{
	colorlinks=true,
	linkcolor=avared,
	urlcolor=avared,
	citecolor=avared,
}

\title{Deploying an L1 on Avalanche}
\subtitle{Workshop · June 28, 2026}
\author{Pedro Coelho do Nascimento}
\institute{Team1}
\date{}

% ══════════════════════════════════════════════════════════════════════════════
\begin{document}
% ══════════════════════════════════════════════════════════════════════════════

\begin{frame}[plain]
	\titlepage
\end{frame}

% ── Who Am I ──────────────────────────────────────────────────────────────────
\begin{frame}{Who Am I}
	\begin{columns}[c]
		\begin{column}{0.9\textwidth}
			{\large\textbf{Pedro Coelho do Nascimento}}\\[0.2em]
			{\small\textcolor{avagraydark}{Rust Back-end Engineer · Full-Stack \& Blockchain Developer · São Paulo, Brazil}}

			\vspace{0.8em}
			\begin{itemize}
				\item Rust engineer at \textbf{Bipa}: payment infrastructure, Central Bank compliance, AML
				\item Web3 open-source contributor: \textbf{Starknet}, \textbf{Stellar} ecosystems (OnlyDust)
				\item Developer at \textbf{Web3EduBrasil}: on-chain certificate issuance platform
			\end{itemize}

			\vspace{0.8em}
			{\small Representing \textbf{Team1}}
		\end{column}
	\end{columns}
\end{frame}

% ── Agenda ────────────────────────────────────────────────────────────────────
\begin{frame}{What We'll Cover}
	\begin{columns}[t]
		\begin{column}{0.48\textwidth}
			\begin{enumerate}
				\item Avalanche architecture overview
				\item The new L1 model (ACP-77)
			\end{enumerate}
		\end{column}
		\begin{column}{0.48\textwidth}
			\begin{enumerate}
				\setcounter{enumi}{2}
				\item Hands-on: deploying an L1
				\item What's next?
			\end{enumerate}
		\end{column}
	\end{columns}
\end{frame}

% ── Why an L1? ────────────────────────────────────────────────────────────────
\begin{frame}{Why Build an Avalanche L1?}
	\framesubtitle{Product and technical benefits}

	\begin{columns}[t]
		\begin{column}{0.32\textwidth}
			\avabox{Your own economy}{
				\begin{itemize}
					\item Custom gas token: users pay in your token, not AVAX
					\item Custom fee model and gas parameters
					\item Fee revenue stays in your ecosystem
				\end{itemize}
			}
		\end{column}
		\begin{column}{0.32\textwidth}
			\avabox{Sovereign by design}{
				\begin{itemize}
					\item Your validator set: permissioned or permissionless
					\item Compliance-friendly: no C-Chain exposure
					\item Isolated execution: other chains cannot affect your throughput
				\end{itemize}
			}
		\end{column}
		\begin{column}{0.32\textwidth}
			\avabox{EVM and beyond}{
				\begin{itemize}
					\item Full EVM compatibility: port existing contracts
					\item Precompiles for minting, access control, fees
					\item ICM for native cross-chain messaging
				\end{itemize}
			}
		\end{column}
	\end{columns}
\end{frame}

\begin{frame}{Who Is It For?}
	\framesubtitle{Common use cases}

	\begin{columns}[t]
		\begin{column}{0.23\textwidth}
			\avabox{Enterprise}{
				Permissioned validators, private chain, regulatory compliance.
			}
		\end{column}
		\begin{column}{0.23\textwidth}
			\avabox{Gaming}{
				Fast finality, gas-free UX for players, custom reward token.
			}
		\end{column}
		\begin{column}{0.23\textwidth}
			\avabox{DeFi}{
				Isolated liquidity, predictable costs, \textasciitilde1--2s settlement.
			}
		\end{column}
		\begin{column}{0.23\textwidth}
			\avabox{Payments}{
				High throughput, stable fees, custom settlement token.
			}
		\end{column}
	\end{columns}
\end{frame}

% ══════════════════════════════════════════════════════════════════════════════
\section{Avalanche Architecture}
% ══════════════════════════════════════════════════════════════════════════════

\begin{frame}{The Primary Network}
	The \textbf{Primary Network} is the foundational layer — every validator on Avalanche
	must validate it.

	\vspace{0.6em}
	\begin{columns}[t]
		\begin{column}{0.32\textwidth}
			\avabox{P-Chain}{
				\textbf{Platform Chain}\\[0.3em]
				Coordinates validators, tracks active L1s, manages staking.
			}
		\end{column}
		\begin{column}{0.32\textwidth}
			\avabox{C-Chain}{
				\textbf{Contract Chain}\\[0.3em]
				EVM-compatible. Home of DeFi, smart contracts, AVAX token.
			}
		\end{column}
		\begin{column}{0.32\textwidth}
			\avabox{X-Chain}{
				\textbf{Exchange Chain}\\[0.3em]
				UTXO-based. Optimized for fast asset creation and transfers.
			}
		\end{column}
	\end{columns}
\end{frame}

\begin{frame}[shrink=5]{Snowman Consensus}
	\framesubtitle{Sub-sampled voting for fast finality}

	\begin{columns}[c]
		\begin{column}{0.55\textwidth}
			\begin{itemize}
				\item Each validator repeatedly samples a small random subset of peers
				\item Switches preference when an $\alpha_p$-majority agrees on a different option
				\item Increments confidence counter when an $\alpha_c$-majority confirms the same option
				\item Decision finalized after $\beta$ consecutive confirming rounds
			\end{itemize}
		\end{column}
		\begin{column}{0.4\textwidth}
			\avabox{Safety guarantee}{
				$$P < \!\left(1 - \frac{\alpha_c}{k}\right)^{\!\beta} \approx 10^{-12}$$
				Liveness requires ${>}75\%$ honest stake.
			}
		\end{column}
	\end{columns}

	\vspace{0.3em}
	{\small\begin{tabular}{@{}lll@{}}
		\toprule
		Symbol & Name & Default \\
		\midrule
		$k$          & Sample size (peers queried per round)            & 20 \\
		$\alpha_p$   & AlphaPreference (votes to switch preference)     & 15 \\
		$\alpha_c$   & AlphaConfidence (votes to increment counter)     & 15 \\
		$\beta$      & Decision threshold (consecutive rounds)          & 20 \\
		\bottomrule
	\end{tabular}}
\end{frame}

\begin{frame}{Avalanche L1s}
	\framesubtitle{Independent blockchains with custom rules}

	\begin{block}{Definition}
		An Avalanche L1 is an independent blockchain with its own validator set,
		tokenomics, and execution layer. Validators come to consensus and validate
		only that blockchain.
	\end{block}

	\vspace{0.5em}
	\begin{columns}[t]
		\begin{column}{0.48\textwidth}
			\textbf{Each L1 specifies its own:}
			\begin{itemize}
				\item Execution logic (VM)
				\item Fee structure
				\item State and networking
				\item Security model
			\end{itemize}
		\end{column}
		\begin{column}{0.48\textwidth}
			\textbf{Network-level benefits:}
			\begin{itemize}
				\item Independent execution threads
				\item Higher aggregate throughput
				\item Isolated state and storage
				\item Cross-chain via ICM
			\end{itemize}
		\end{column}
	\end{columns}
\end{frame}

% ══════════════════════════════════════════════════════════════════════════════
\section{ACP-77: Reinventing Subnets}
% ══════════════════════════════════════════════════════════════════════════════

\begin{frame}{The Old Model: Subnets}
	\framesubtitle{Before Avalanche9000}

	\begin{columns}[c]
		\begin{column}{0.55\textwidth}
			\begin{itemize}
				\item Subnet validators \textbf{required} to also validate the Primary Network
				\item Mandatory \textbf{2,000 AVAX} stake per validator on Primary Network
				\item At \$20/AVAX that's \textbf{\$40,000 locked per validator}
				\item Minimum viable Subnet needed 5 validators = \$200,000
				\item Validator set managed directly on the P-Chain by the Subnet owner
			\end{itemize}
		\end{column}
		\begin{column}{0.4\textwidth}
			\avawarning{Bottleneck}{
				Every Subnet validator had to run and sync the full Primary Network,
				increasing hardware costs and coupling performance to Primary Network load.
			}
		\end{column}
	\end{columns}
\end{frame}

\begin{frame}[shrink=5]{Avalanche9000: The Etna Upgrade}
	\framesubtitle{The largest network upgrade in Avalanche's history}

	\begin{columns}[t]
		\begin{column}{0.48\textwidth}
			\avabox{Economic barriers removed}{
				2,000 AVAX stake requirement \textbf{eliminated}. L1 validators pay only a
				continuous P-Chain fee of \textbf{1--10 AVAX/month} per validator.
			}
			\vspace{0.15em}
			\avabox{Running costs reduced}{
				Validators only run the L1, no need to sync the Primary Network.
				Lower hardware requirements, simpler ops.
			}
		\end{column}
		\begin{column}{0.48\textwidth}
			\avabox{Improved compliance}{
				Validator sets are fully separate from the Primary Network.
				Institutions can launch private L1s with \textbf{zero exposure} to public
				permissionless chains.
			}
			\vspace{0.15em}
			\avabox{Fault isolation}{
				Congestion or outages on the Primary Network no longer affect L1
				performance, and vice versa.
			}
		\end{column}
	\end{columns}

	\vspace{0.2em}
	\begin{center}
		\small\textcolor{avagraydark}{Implemented via ACP-77 (Reinventing Subnets), supersedes ACP-13}
	\end{center}
\end{frame}

\begin{frame}{ACP-77: What Changed}
	\framesubtitle{Activated · Standards Track}

	\begin{columns}[t]
		\begin{column}{0.32\textwidth}
			\avabox{1. Separation}{
				Subnet validators no longer required to validate the Primary Network.
				2,000 AVAX requirement removed.
			}
		\end{column}
		\begin{column}{0.32\textwidth}
			\avabox{2. Sovereignty}{
				Validator set management moves from P-Chain to a
				\textbf{Validator Manager contract} on the L1 itself.
			}
		\end{column}
		\begin{column}{0.32\textwidth}
			\avabox{3. Continuous fees}{
				L1 validators pay a continuous fee to the P-Chain — no large upfront
				stake needed.
			}
		\end{column}
	\end{columns}

	\vspace{0.5em}
	\begin{center}
		\small\textcolor{avagraydark}{ACP-77 supersedes ACP-13 (Subnet-Only Validators)}
	\end{center}
\end{frame}

\begin{frame}{Subnet vs.\ L1: Side by Side}
	\begin{center}
	\renewcommand{\arraystretch}{1.4}
	\begin{tabular}{@{} l >{\raggedright}p{4.5cm} >{\raggedright\arraybackslash}p{4.5cm} @{}}
		\toprule
		& \textbf{Subnet (legacy)} & \textbf{L1 (ACP-77)} \\
		\midrule
		Validator set    & Must validate Primary Network & Sovereign, custom set \\
		Stake required   & 2,000 AVAX minimum            & No minimum (continuous fee) \\
		Set management   & P-Chain / subnet owner        & Validator Manager contract \\
		Primary net sync & Full sync required            & Header-only (partial sync) \\
		Finality         & ~1–2 s                        & ~1–2 s \\
		Interoperability & ICM                           & ICM \\
		\bottomrule
	\end{tabular}
	\end{center}
\end{frame}

% ══════════════════════════════════════════════════════════════════════════════
\section{Deploying Your L1}
% ══════════════════════════════════════════════════════════════════════════════

\begin{frame}{What We're Building}
	\framesubtitle{Four-step process}

	\begin{center}
	\begin{tikzpicture}[
		box/.style={
			rectangle, rounded corners=2pt,
			draw=avared, fill=avagray,
			text width=2.6cm, align=center,
			minimum height=1.2cm, font=\small\sffamily
		},
		arr/.style={->, thick, color=avared, >=latex}
	]
		\node[box] at (0,0)   (n1) {\textbf{1.} Create Subnet \& Chain};
		\node[box] at (3.2,0) (n2) {\textbf{2.} Run a Validator Node};
		\node[box] at (6.4,0) (n3) {\textbf{3.} Convert to L1};
		\node[box] at (9.6,0) (n4) {\textbf{4.} Test \& Deploy};

		\draw[arr] (n1.east) -- (n2.west);
		\draw[arr] (n2.east) -- (n3.west);
		\draw[arr] (n3.east) -- (n4.west);
	\end{tikzpicture}
	\end{center}

	\vspace{0.6em}
	\begin{columns}[t]
		\begin{column}{0.48\textwidth}
			\textbf{Prerequisites:}
			\begin{itemize}
				\item Core Wallet (set to Testnet)
				\item Testnet AVAX on P-Chain
				\item AWS EC2 instance (we'll set it up together)
			\end{itemize}
		\end{column}
		\begin{column}{0.48\textwidth}
			\textbf{Tools:}
			\begin{itemize}
				\item \href{https://build.avax.network/tools/l1-toolbox}{Avalanche Builder Console}
				\item Docker on EC2 for AvalancheGo
				\item MetaMask / Core Wallet for RPC testing
			\end{itemize}
		\end{column}
	\end{columns}
\end{frame}

\begin{frame}[shrink=5]{Node Roles}
	\framesubtitle{Validator · RPC · Combined}

	\begin{columns}[t]
		\begin{column}{0.32\textwidth}
			\avabox{Validator node}{
				\textbf{Port 9651} (P2P)\\[0.3em]
				Participates in consensus and produces blocks.
				RPC bound to \texttt{localhost} only. Not exposed publicly.
			}
		\end{column}
		\begin{column}{0.32\textwidth}
			\avabox{RPC node}{
				\textbf{Port 9650} (HTTP)\\[0.3em]
				Serves JSON-RPC to wallets and dApps.
				Not part of the validator set. No consensus role.
			}
		\end{column}
		\begin{column}{0.32\textwidth}
			\avabox{Validator + RPC}{
				\textbf{Ports 9650 + 9651}\\[0.3em]
				Both roles on one machine.
				Simpler ops, higher resource use.
				\textbf{What we use today.}
			}
		\end{column}
	\end{columns}

	\vspace{0.4em}
	\avabox{Workshop topology}{
		Single EC2 running AvalancheGo in \textbf{validator-rpc} mode.
		\textbf{9651}: open inbound/outbound for P2P consensus.
		\textbf{9650}: bound to localhost; a reverse proxy (nginx / Caddy) on \textbf{443/80} terminates TLS and forwards to it.
		Adequate for testnet. Production separates validator and RPC nodes.
	}
\end{frame}

% ══════════════════════════════════════════════════════════════════════════════
\section{What's Next?}
% ══════════════════════════════════════════════════════════════════════════════

\begin{frame}{Customizing Your L1}
	\framesubtitle{EVM precompiles — Solidity contracts at known addresses}

	\begin{columns}[t]
		\begin{column}{0.48\textwidth}
			\begin{itemize}
				\item \textbf{NativeMinter} — mint your gas token on-chain
				\item \textbf{FeeManager} — adjust fee config without a hard fork
				\item \textbf{RewardManager} — configure block reward distribution
			\end{itemize}
		\end{column}
		\begin{column}{0.48\textwidth}
			\begin{itemize}
				\item \textbf{ContractDeployerAllowList} — restrict who can deploy contracts
				\item \textbf{TxAllowList} — restrict who can send transactions
				\item \textbf{WarpMessenger} — send cross-chain messages via ICM
			\end{itemize}
		\end{column}
	\end{columns}

	\vspace{0.25em}
	\avabox{Enable in genesis}{
		Precompiles are activated in \texttt{genesis.json} under the
		\texttt{config} block. Access control lists define admin
		and enabled addresses.
	}
\end{frame}

\begin{frame}[shrink=10]{Going Further}
	\framesubtitle{Natural next steps from what we built today}

	\begin{columns}[t]
		\begin{column}{0.48\textwidth}
			\avabox{PoA $\to$ PoS}{
				Replace the PoA manager with a \textbf{Staking Manager}.
				Anyone who stakes your native token (or an ERC-20) can join the
				validator set, no owner approval needed.
			}
		\end{column}
		\begin{column}{0.48\textwidth}
			\avabox{Multisig governance}{
				Wrap the PoA manager in a \textbf{Safe multisig} so no single key
				controls validator adds/removes. Critical before mainnet.
			}
		\end{column}
	\end{columns}

	\vspace{0.2em}
	\begin{columns}[t]
		\begin{column}{0.48\textwidth}
			\avabox{ICM: Interchain Messaging}{
				Send arbitrary messages between Avalanche L1s.
				Token bridges, cross-chain app logic, shared state.
				Needs a \textbf{Warp relayer} running alongside your node.
			}
		\end{column}
		\begin{column}{0.48\textwidth}
			\avabox{Beyond EVM}{
				\textbf{HyperSDK}: build a custom VM in Go with full control
				over state, execution, and fee models. For when the EVM is
				the bottleneck, not the answer.
			}
		\end{column}
	\end{columns}
\end{frame}

\begin{frame}{Resources}
	\begin{columns}[t]
		\begin{column}{0.48\textwidth}
			\textbf{Documentation \& tooling:}
			\begin{itemize}
				\item \href{https://build.avax.network}{build.avax.network} — Builder Hub
				\item \href{https://build.avax.network/tools/l1-toolbox}{L1 Toolbox}
				\item \href{https://build.avax.network/academy}{Academy courses}
			\end{itemize}

			\vspace{0.5em}
			\textbf{Key ACPs:}
			\begin{itemize}
				\item ACP-77: Reinventing Subnets
				\item ACP-13: Subnet-Only Validators (superseded)
			\end{itemize}
		\end{column}
		\begin{column}{0.48\textwidth}
			\textbf{Community:}
			\begin{itemize}
				\item \href{https://discord.gg/avax}{discord.gg/avax}
				\item \href{https://github.com/ava-labs}{github.com/ava-labs}
			\end{itemize}

			\vspace{0.5em}
			\textbf{Today's code:}
			\begin{itemize}
				\item \cli{ava-labs/avalanchego}
				\item \cli{ava-labs/subnet-evm}
				\item \cli{ava-labs/icm-contracts}
			\end{itemize}
		\end{column}
	\end{columns}
\end{frame}

% ── Conclusion ────────────────────────────────────────────────────────────────
\begin{frame}{What We Built}
	\framesubtitle{From zero to a live Avalanche L1 on Fuji}

	\begin{columns}[t]
		\begin{column}{0.48\textwidth}
			\avabox{Your L1: live on Fuji}{
				\begin{itemize}
					\item Custom subnet \& blockchain registered on P-Chain
					\item AvalancheGo validator node running via Docker
					\item Validator Manager contract (PoA) deployed
					\item RPC endpoint testable with MetaMask / Core
				\end{itemize}
			}
		\end{column}
		\begin{column}{0.48\textwidth}
			\avabox{Key concepts covered}{
				\begin{itemize}
					\item Primary Network: P/C/X-Chain separation
					\item ACP-77: sovereign validator sets, no 2,000 AVAX requirement
					\item Validator Manager: on-chain governance
					\item Node roles: validator, RPC, combined topology
				\end{itemize}
			}
		\end{column}
	\end{columns}
\end{frame}

% ── Q&A ───────────────────────────────────────────────────────────────────────
\begin{frame}[plain]
	\begin{tikzpicture}[remember picture, overlay]
		\fill[avadark] (current page.north west) rectangle (current page.south east);
		\draw[avared, line width=3pt]
			(current page.north west) rectangle (current page.south east);
	\end{tikzpicture}
	\vfill
	\begin{center}
		{\usebeamerfont{title}\usebeamercolor[fg]{title}Questions?\par}
		\vspace{1em}
		{\usebeamerfont{subtitle}\usebeamercolor[fg]{subtitle}Team1\par}
	\end{center}
	\vfill
\end{frame}

\end{document}
